\chapter{Setup and Installation}
\label{ch:setup}

This chapter walks through downloading, installing, and running Kiss for the first time.  By the end you will have a working development environment with both the frontend and backend servers running.

\section{Prerequisites}
\label{sec:setup-prerequisites}

Kiss requires:

\begin{itemize}
  \item \textbf{Java 17 or later} (JDK, not just JRE).  Kiss compiles and runs Java and Groovy code at runtime, so a full JDK is required.  Verify with \texttt{java -version} and \texttt{javac -version}.
  \item \textbf{A modern web browser} --- Chrome, Firefox, Edge, or Safari.
  \item \textbf{Optionally, an IDE} for code editing.  See Appendix~\ref{ch:ide-setup} for setup instructions.  Any text editor works; Kiss does not depend on IDE tooling.
\end{itemize}

\begin{notebox}
Kiss includes its own build system, an embedded Tomcat server, and all required JAR dependencies.  No Maven, Gradle, or Node.js installation is necessary.
\end{notebox}

\begin{warnbox}
Windows requires several additional system-level dependencies (such as a package manager and supporting tools) that Linux and macOS typically provide out of the box.  These are not duplicated here --- consult the \emph{Kiss Getting Started Guide} for the complete Windows setup procedure before proceeding.
\end{warnbox}

\section{Downloading Kiss}
\label{sec:setup-download}

Kiss is distributed as a self-contained directory.  Clone the repository or download and extract the archive:

\begin{lstlisting}[style=shell]
git clone https://github.com/blakemcbride/Kiss.git
cd Kiss
\end{lstlisting}

\noindent
Alternatively, download a release archive from the project's GitHub page, extract it, and \texttt{cd} into the resulting directory.

\section{Directory Structure}
\label{sec:setup-directory}

The top-level directory contains everything needed to build and run the application:

\begin{lstlisting}
Kiss/
  src/
    main/
      core/          # Core framework (do not modify)
      backend/       # Your application services
      frontend/      # Your application screens
        kiss/        # Framework components (do not modify)
      precompiled/   # Shared Java utilities
    test/            # Unit tests
  libs/              # Third-party JAR dependencies
  manual/            # Documentation source
  tomcat/            # Embedded Tomcat server
    logs/            # Server log files
  work/              # Build output
\end{lstlisting}

\begin{warnbox}
Two directories must not be modified:
\begin{itemize}
  \item \texttt{src/main/core/} --- the Kiss framework's Java classes.
  \item \texttt{src/main/frontend/kiss/} --- the Kiss frontend component library.
\end{itemize}
Your application code goes in \texttt{src/main/backend/} (services) and \texttt{src/main/frontend/} (screens, excluding the \texttt{kiss/} subdirectory).
\end{warnbox}

\section{First Run}
\label{sec:setup-first-run}

Start both the frontend and backend development servers with a single command:

\begin{lstlisting}[style=shell]
./bld -v develop              [Linux, macOS, BSD, etc.]
      -or-
bld -v develop                [Windows]
\end{lstlisting}

\begin{warnbox}
On Linux, macOS, BSD, and other Unix-like systems, run the script as \texttt{./bld} (the leading \texttt{./} is required).  On Windows, run it as \texttt{bld} (without the \texttt{./} prefix), which invokes the \texttt{bld.cmd} batch file.  All examples in this book use the \texttt{./bld} form; Windows users should drop the \texttt{./} prefix in every command.
\end{warnbox}

\noindent
This command:
\begin{enumerate}
  \item Compiles the core framework and any precompiled utilities.
  \item Starts the backend server on port 8001.
  \item Starts a frontend development server on port 8000.
  \item Watches for file changes and recompiles automatically.
\end{enumerate}

\noindent
Open \texttt{http://localhost:8000} in a browser.  The login page should appear.

\begin{tipbox}
Keep the terminal running.  You can edit backend and frontend files while the system is running, and changes will take effect without a restart.
\end{tipbox}

\section{Verifying the Installation}
\label{sec:setup-verify}

Log in with the default credentials: username \texttt{kiss}, password \texttt{password}.  These credentials are valid when using the bundled SQLite demo database.

After logging in, the example application's navigation menu appears.  Click through the screens to confirm that the database, frontend, and backend are communicating correctly:

\begin{description}[style=nextline]
  \item[Controls] A showcase of every Kiss UI component: text inputs, dropdowns, date pickers, checkboxes, radio buttons, and more.
  \item[CRUD] A phone book application demonstrating create, read, update, and delete operations with a data grid, popup editor, PDF report, and CSV export.
  \item[REST Services] Examples of calling backend service methods and displaying results.
  \item[Users] User management for the built-in authentication system.
  \item[File Upload] Demonstrates file upload handling.
\end{description}

\section{The Example Application}
\label{sec:setup-example-app}

Kiss ships with a working example application that serves two purposes: it verifies that everything is installed correctly, and it provides a template for your own code.  The example includes:

\begin{itemize}
  \item A login page with session management.
  \item A navigation framework that loads screens dynamically.
  \item A CRUD screen (phone book) with a grid, popup editor, PDF report, and CSV export.
  \item A controls showcase demonstrating every UI component.
  \item A file upload screen.
  \item An AI/LLM integration screen (requires Ollama).
\end{itemize}

\noindent
All example services are in \texttt{src/main/backend/services/} and all example screens are in \texttt{src/main/frontend/screens/}.  Study these files as you read through the book --- they are the simplest working examples of each feature.

\section{Development URLs}
\label{sec:setup-urls}

During development, two servers run simultaneously:

\begin{longtable}{ll}
\toprule
\textbf{URL} & \textbf{Purpose} \\
\midrule
\texttt{http://localhost:8000} & Frontend (HTML/CSS/JS) \\
\texttt{http://localhost:8001} & Backend (JSON-RPC services) \\
\bottomrule
\end{longtable}

\noindent
Both ports come from a single \texttt{portBase} variable in \texttt{Tasks.java} (front end = \texttt{portBase}, back end = \texttt{portBase}~+~1; default 8000), so several Kiss applications can run in development at once by giving each a different base --- see Section~\ref{sec:build-ports}.  Pages served by the frontend dev server call the backend at their own port~+~1 automatically.  In production, both are served from the same Tomcat instance (see Chapter~\ref{ch:deployment}).

\section{Server Logs}
\label{sec:setup-logs}

Backend log output is written to \texttt{tomcat/logs/catalina.out}.  This file contains startup messages, service compilation events, request logs, and any errors thrown by backend code.

On Linux, macOS, BSD, and other Unix-like systems, use the bundled \texttt{view-log} script to tail the log:

\begin{lstlisting}[style=shell]
./view-log
\end{lstlisting}

\noindent
On Windows, the log window opens automatically when the backend starts --- no separate command is needed.

This is the first place to look when something goes wrong.
